% !TEX root = ../thuthesis-example.tex
\chapter{研究成果}

\section{暗物质场单点统计}
    我们比较新流程与ZA基准分别与N体模拟的过密度 \(\delta(\boldsymbol{x})=\rho(\boldsymbol{x})/\bar{\rho}-1\) 的残差的概率分布函数（PDF），即比较 \(\mathrm{PDF}(\delta_{\text{new}} - \delta_{\text{N-body}})\) 和 \(\mathrm{PDF}(\delta_{\text{ZA}} - \delta_{\text{N-body}})\)。其中 \(\rho(\boldsymbol{x})\) 是位置 \(\boldsymbol{x}\) 处的暗物质密度，\(\bar{\rho}\) 是全盒子的平均暗物质密度；\(\delta_{\text{new}}\)、\(\delta_{\text{ZA}}\) 和 \(\delta_{\text{N-body}}\) 分别表示新流程、ZA 基准和 N 体模拟得到的过密度场。
结果显示，新流程与N体模拟的过密度残差分布更接近0且更窄，表明新流程在单点统计上相较于ZA基准有明显改进。同时，我们也绘制了密度场的视觉直观展示，结果显示新流程的密度场与N体模拟更为相似，在小尺度上相比于，ZA基线，神经网络能使粒子更加紧凑，细节更加丰富，
进一步验证了新流程在单点统计上的改进。
（TODO:realization 32的对比图两张：第一张是density plot，包括Nbody的、ZA的、ALPT的和新流程的；第二张是残差PDF和Nbody的对比，包括ZA的、ALPT的和新流程的）

\section{暗物质场两点统计}
    我们比较新流程与多种微扰论方法与N体模拟的物质功率谱$P(k)$比值，物质功率谱的定义为
\begin{equation}
    P(k) = \langle |\delta(\boldsymbol{k})|^2 \rangle,
\end{equation}
其中 \(\delta(\boldsymbol{k})\) 是过密度场 \(\delta(\boldsymbol{x})\) 的傅里叶变换，\(\boldsymbol{k}\) 是波矢，\(k=|\boldsymbol{k}|\) 是波数；\(\langle\cdot\rangle\) 表示对同一 \(k\) 壳层内所有傅里叶模式取平均。结果显示，新流程一直到Nyquist频率（灰色虚线）都与N体模拟的功率谱比值接近1，表明新流程在两点统计上相较于其他方法有明显改进，
几乎对N体模拟进行了完全恢复。Nyquist频率以上的波数处于欠采样状态，更高的$k$处的结果无物理意义。
（TODO:32-47共16个模拟的功率谱图，分为两个子图，上下排列，上面是功率谱绝对值图，包括Nbody的、ZA的、ALPT的和新流程的；下面是功率谱比值图，包括ZA的、ALPT的和新流程的，用
不同颜色绘制，绘制平均值和误差棒，k最大画到1，还要画Nyquist频率；下面是与Nbody的比值图，同样绘制平均值和误差棒，k最大画到1，还要画Nyquist频率）

\section{暗物质场三点统计}
    我们比较新流程与多种微扰论方法与N体模拟的物质双谱$B(k_1, k_2, k_3)$比值，物质双谱的定义为
\begin{equation}
    B(k_1, k_2, k_3) = \langle \delta(\boldsymbol{k}_1) \delta(\boldsymbol{k}_2) \delta(\boldsymbol{k}_3) \rangle,
\end{equation}
其中三个波矢满足构成矢量三角形的条件 \(\boldsymbol{k}_1 + \boldsymbol{k}_2 + \boldsymbol{k}_3 = 0\)，\(k_1, k_2, k_3\) 分别是三个波矢的模；\(\langle\cdot\rangle\) 表示对满足该三角形构型的傅里叶模式组合取平均。在比较时，我们固定其中两个波数 \(k_1\) 和 \(k_2\)，并将两波矢夹角 \(\theta\) 作为横轴。
我们选择特定几个配置进行比较，结果显示，在所有配置下，新流程与N体模拟的双谱比值都接近1，表明新流程在三点统计上相较于其他方法有明显改进，几乎对N体模拟进行了完全恢复。
（TODO:32-47共16个模拟的双谱比值图，包括ZA的、ALPT的和新流程的，用不同颜色绘制，绘制平均值和误差棒，k最大画到1，还要画Nyquist频率。
不同方法分别用不同线型表示，新流程用实线、ZA用虚线、ALPT用点线。选用不同configuration绘制到同一张图：分别用:k1=0.05,k2=0.1;k1=0.1,k2=0.2;k1=0.2,k2=0.4;
k1=0.1,k2=0.1;k1=0.2,k2=0.2;k1=0.4,k2=0.4.）

\section{暗物质功率谱协方差矩阵}
    在这里我们比较了48个模拟得到的暗物质功率谱的协方差矩阵。结果显示，新流程得到的功率谱协方差矩阵与N体模拟的协方差矩阵在结构上非常相似，表明新流程在统计性质上相较于其他方法有明显改进，
但在小尺度上协方差对角元过小，表明新流程在小尺度上对N体模拟的功率谱恢复过度，导致了欠估计的误差，但目前用于计算协方差矩阵的
样本数量过小，结果可能不稳定，未来需要更多的模拟来验证这一结果。
（TODO:16-63共4个模拟的功率谱协方差矩阵图，包括Nbody的、ZA的、ALPT的和新流程的，用linear bin；第二张图是不同方法的协方差矩阵对角元与Nbody比值图，k最大画到1，linear bin，还要绘制Nyquist频率）

\section{速度场视觉展示}
    我们比较新流程与ZA基准与N体模拟的速度场的视觉展示。结果显示，新流程得到的速度场在视觉上更接近N体模拟，尤其在小尺度上相比于ZA基线，
神经网络能使粒子具有更丰富的速度结构，进一步验证了新流程在速度场恢复上的改进。
（TODO:realization 32的速度场对比图，包括Nbody的、ZA的和新流程的,以及Nbody-新流程的残差）

\section{速度场两点统计}
（TODO：绘制速度场散度功率谱、旋度功率谱以及分量功率谱与Nbody的比值图，包括ZA的和新流程的，用不同颜色绘制，绘制平均值和误差棒，k最大画到1，还要画Nyquist频率）


